\section{Skills 的分发与管理}

\subsection{两种分发方式}

Skills 最大的好处之一是可以与团队其他成员共享。有两种方式：

\begin{enumerate}
    \item \textbf{提交到代码仓库}：把 Skills 放在 \texttt{./.claude/skills} 下。适合在较少代码仓库上协作的小团队，但每个被 check in 的 Skill 都会给模型上下文增加一点负担。
    \item \textbf{内部插件市场}：搭建 Claude Code Plugin Marketplace，让用户可以上传和安装插件。随着规模扩大，内部插件市场让团队成员自己决定安装哪些 Skills，避免上下文膨胀。
\end{enumerate}

\subsection{管理插件市场（Managing a Marketplace）}

怎么决定哪些 Skills 进入市场？人们怎么提交？

Anthropic 内部\textbf{没有一个中心团队}来做决定，而是让最有用的 Skills 有机地（organically）涌现出来：

\begin{enumerate}
    \item 如果你有一个想让别人试试的 Skill，先上传到 GitHub 的 sandbox 文件夹，然后在 Slack 等论坛推广
    \item 当 Skill 获得了足够的关注（由 Skill 作者自己判断），就可以提交 PR 将其移到正式的 marketplace
    \item 在正式发布前确保有审核机制
\end{enumerate}

\begin{warningbox}{质量控制不可缺少}
创建质量差或重复的 Skills 非常容易（``it can be quite easy to create bad or redundant skills''）。在正式发布前确保有某种审核机制（curation method）非常重要。
\end{warningbox}

\subsection{组合 Skills（Composing Skills）}

Skills 之间可以互相依赖。例如，你可能有一个文件上传 Skill 负责上传文件，还有一个 CSV 生成 Skill 负责生成 CSV 并调用上传。这种依赖管理目前还\textbf{没有被原生支持}（not natively built into marketplaces or skills yet），但你可以直接按名字引用其他 Skills --- 只要对方已安装，模型就会调用它们。

\subsection{衡量 Skills 的效果（Measuring Skills）}

\begin{importantbox}{用 PreToolUse 钩子追踪 Skill 使用情况}
Anthropic 使用一个 \texttt{PreToolUse} 钩子在公司内部记录 Skill 的使用情况。这让他们能够发现哪些 Skills 受欢迎，以及哪些 Skills 的触发频率低于预期（undertriggering）。这是理解 Skill 生态健康状况的关键手段。
\end{importantbox}

\subsection{本章小结}

Skills 的分发有仓库级和市场级两种方式，适合不同规模的团队。管理的核心思路是让好的 Skills 自然涌现而非自上而下指定，同时做好质量控制。Skills 之间可以按名字互相引用实现组合，但原生依赖管理尚待完善。通过 PreToolUse 钩子追踪使用数据是衡量 Skill 生态健康的有效手段。

